\clearpage
\VerdierPhysicalPageAnchor{0163}%
\begingroup\small\setlength{\parskip}{0pt}\setlength{\emergencystretch}{1.5em}
\VerdierSemanticAnchor{ch4-4-1-2-b}
b) A functor from $\mathcal{F}\ell_2(\mathcal J)$ to the category of
distinguished triangles in $\mathcal D$, of the form
\[
\omega_2=(\mathord\bullet\xrightarrow f\mathord\bullet\xrightarrow g\mathord\bullet)
\longmapsto
X_f\xrightarrow{X*u_1(\omega_2)}X_{gf}
\xrightarrow{X*u_2(\omega_2)}X_g
\xrightarrow{\delta(\omega_2)}X_f[1].
\]
To denote spectral objects of type $\mathcal J$, we shall use the notation
\[
(X_f, f\in\operatorname{Ob}\mathcal{F}\ell_1(\mathcal J),\
\delta(\omega_2),
\ \omega_2\in\operatorname{Ob}\mathcal{F}\ell_2(\mathcal J)),
\]
or simply $(X_f,\delta)$ when no confusion can result.

\VerdierSemanticAnchor{ch4-4-1-3}
\medskip\noindent\textbf{4.1.3.}\quad
Let $C$ be an additive category and let
$F\colon\mathcal J\longrightarrow\operatorname{comp}(C)$ be a functor. In
Chapter I, (3.3.6), we showed how to associate with $F$ a spectral object of
type $\mathcal J$ with values in $K(C)$. The spectral object described in
Chapter I, (3.3.6), is called the \emph{spectral object associated with the
functor $F$}.

\VerdierSemanticAnchor{ch4-4-1-4}
\medskip\noindent\textbf{4.1.4.}\quad
Let $\mathcal A$ be an abelian category and $\mathcal J$ a category. A
\emph{spectral object of type $\mathcal J$ with values in $\mathcal A$}
consists of:

a) A sequence of functors
$H^n\colon\mathcal{F}\ell_1(\mathcal J)\longrightarrow\mathcal A$,
$n\in\mathbb Z$.

b) For every object
$\omega_2=(\mathord\bullet\xrightarrow f\mathord\bullet\xrightarrow g\mathord\bullet)$
of $\mathcal{F}\ell_2(\mathcal J)$ and every integer $n$, a morphism
$\delta^n(\omega_2)\colon H^n(g)\to H^{n+1}(f)$, functorial in $\omega_2$,
such that the diagram
\[
\cdots\longrightarrow H^n(f)
\xrightarrow{H^n*u_1(\omega_2)}H^n(gf)
\xrightarrow{H^n*u_2(\omega_2)}H^n(g)
\xrightarrow{\delta^n(\omega_2)}H^{n+1}(f)
\longrightarrow\cdots
\]
is an infinite exact sequence in $\mathcal A$.

We shall use the notation
\[
(H^n(f), f\in\operatorname{Ob}\mathcal{F}\ell_1(\mathcal J),\
\delta^n(\omega_2),
\ \omega_2\in\operatorname{Ob}\mathcal{F}\ell_2(\mathcal J))
\]
to denote spectral objects of type $\mathcal J$, or simply
$(H^n(f),\delta)$ when no confusion can result.

\VerdierSemanticAnchor{ch4-4-1-5}
\medskip\noindent\textbf{4.1.5.}\quad
Let $F\colon\mathcal D\to\mathcal D'$ be an exact functor between two
triangulated categories. Let $(X_f,\delta)$ be a spectral object of type
$\mathcal J$ with values in $\mathcal D$. The pair of functors
$(FX_f,F\delta)$ is a spectral object of type $\mathcal J$ with values in
$\mathcal D'$. Similarly, a cohomological functor
$H\colon\mathcal D\to\mathcal A$ from a triangulated category to an
abelian category transforms spectral objects with values in $\mathcal D$
into spectral objects with values in $\mathcal A$.

\VerdierSemanticAnchor{ch4-4-1-6}
\medskip\noindent\textbf{4.1.6.}\quad
Let $\mathcal D$ be a triangulated category and let $(X_f,\delta)$,
$(Y_f,\delta')$ be two spectral objects of type $\mathcal J$ with values in
$\mathcal D$. A \emph{morphism of spectral objects} is a morphism of
functors
\[
m_f\colon X_f\longrightarrow Y_f,\qquad
f\in\operatorname{Ob}\mathcal{F}\ell_1(\mathcal J),
\]
\endgroup
